\documentclass[11pt,a4paper]{article}
\usepackage[T1]{fontenc}
\usepackage{lmodern}
\usepackage[margin=25mm]{geometry}
\usepackage{amsmath,amssymb,amsthm,mathtools}
\usepackage{enumitem,microtype,xcolor,booktabs,tabularx}
\usepackage[hyperpageref]{backref}
\usepackage[pagebackref]{hyperref}
\hypersetup{
    colorlinks=true,
    linkcolor=red,
    filecolor=blue,
    urlcolor=orange,
    citecolor=violet
}
\renewcommand*{\backref}[1]{}
\renewcommand*{\backrefalt}[4]{
  \ifcase #1 Not cited.
  \or Cited on page #2.
  \else Cited on pages #2.
  \fi}
\definecolor{darkblue}{rgb}{0,0,0.7}
\newcommand{\defn}[1]{\textsl{\color{darkblue} #1}}
\setlist[enumerate]{leftmargin=*,itemsep=.5em,topsep=.4em}
\setlength{\emergencystretch}{2em}
\newcommand{\D}{\mathcal D}
\renewcommand{\H}{\mathcal H}
\newcommand{\ZZ}{\mathbb{Z}}
\newcommand{\op}{\mathrm{op}}
\newcommand{\RR}{\mathbf{R}}
\newcommand{\std}{\Delta}
\newcommand{\costd}{\nabla}
\newcommand{\dual}{\mathbb D}
\newcommand{\Ker}{\operatorname{Ker}}
\DeclareMathOperator{\fd}{fd}
\DeclareMathOperator{\Hom}{Hom}
\DeclareMathOperator{\End}{End}
\DeclareMathOperator{\REnd}{\RR End}
\DeclareMathOperator{\add}{add}
\DeclareMathOperator{\thick}{thick}
\DeclareMathOperator{\per}{per}
\DeclareMathOperator{\Filt}{Filt}
\DeclareMathOperator{\rad}{rad}
\let\mod\relax
\DeclareMathOperator{\mod}{mod}
\theoremstyle{plain}
\newtheorem{theorem}{Theorem}[section]
\newtheorem{proposition}[theorem]{Proposition}
\newtheorem{lemma}[theorem]{Lemma}
\newtheorem{corollary}[theorem]{Corollary}
\theoremstyle{definition}
\newtheorem{definition}[theorem]{Definition}
\newtheorem{example}[theorem]{Example}
\newtheorem{remark}[theorem]{Remark}
\numberwithin{equation}{section}
\title{$m$-qh dgas: composition factors and shifted filtrations}
\author{}
\date{}


\begin{document}
\maketitle
\begin{abstract}
We define $m$-quasi-hereditary dg algebras by two requirements:
the cohomology of each standard involves only simple modules with
smaller or equal vertex labels in the chosen total order, with its own
simple occurring exactly once in degree zero and never in negative
degrees, and projective presentations have kernels filtered by allowed
shifts of standards indexed by strictly larger vertices in that order.
Exceptionality, fullness, and homological smoothness
are consequences. We prove the equivalence with the exceptional-sequence
formulation, show that the ordered projectives determine the standards
and their left-dual sequence, and establish filtration criteria and
the least possible value of $m$. Graded path
algebras provide examples with arbitrary vertex orders and with
non-trivial differentials. Counterexamples distinguish indispensable
axioms, forbidden orders, insufficient values of $m$, and properties of
classical filtrations that do not survive.
\end{abstract}

\section{The two axioms}\label{sec:axioms}

Fix an algebraically closed field $\Bbbk$ and integers $m,n\geq1$.
All dg modules are right modules; $\D(A)$ denotes their derived category.
Complexes are cochain complexes:
$d_X:X^q\to X^{q+1}$ and $H^q(X[s])=H^{q+s}(X)$.
Let $A$ be a proper connective dg algebra which is $m$-truncated:
\begin{equation}\label{eq:standing}
 \dim_{\Bbbk}H^*(A)<\infty,\qquad
 H^q(A)=0\quad(q\notin[1-m,0]).
\end{equation}
Assume that $B=H^0(A)$ is basic and that closed degree-zero orthogonal
idempotents $e_1,\ldots,e_n$ sum to $1$, with classes forming a complete
set of primitive idempotents of $B$. Write $e_v$ also for its class in
$B$, and let $S_v$ be the simple top of $e_vB$. Set
\begin{equation}\label{eq:categories}
 \begin{gathered}
 P_v=e_vA,\qquad
 \D_{\fd}(A)=\{X\in\D(A):\dim_{\Bbbk}H^*(X)<\infty\},\\
 \per A=\thick_{\D(A)}(A),\qquad
 \H_A^m=\D_{\fd}(A)^{[1-m,0]}.
 \end{gathered}
\end{equation}
Here $\thick$ denotes closure under shifts, cones, and direct summands;
objects of $\per A$ are called \defn{perfect}. The interval in
\eqref{eq:categories} specifies cohomological degrees.
Write $\add A$ for direct summands of finite direct sums of $A$.
For a finite-dimensional $B$-module $M$, write $[M:S_w]$ for its
composition multiplicity. Since $B$ is split basic,
\begin{equation}\label{eq:projective-cohomology}
 [M:S_w]=\dim_{\Bbbk}Me_w,\qquad
 \Hom_{\D_{\fd}(A)}(P_w,X[t])\cong H^t(X)e_w
 \quad(X\in\D_{\fd}(A),\ t\in\ZZ).
\end{equation}

Choose a permutation $\sigma\in S_n$ of the vertex labels. The vertex
at position $a$ is $\sigma(a)$; the position of vertex $v$ is
$\sigma^{-1}(v)$. Define
\begin{equation}\label{eq:order}
 u\leq_\sigma v\quad\Longleftrightarrow\quad
 \sigma^{-1}(u)\leq\sigma^{-1}(v),\qquad
 u<_\sigma v\quad\Longleftrightarrow\quad
 u\leq_\sigma v\text{ and }u\neq v.
\end{equation}
The least and greatest vertices with respect to $\leq_\sigma$ are
$\sigma(1)$ and $\sigma(n)$.
Vertex labels and positions in a sequence are distinct.

The \defn{$m$-extended module category} $\H_A^m$ is extriangulated:
its conflations are the triangles
$X\to Y\to Z\to X[1]$ with $X,Y,Z\in\H_A^m$.
We write $X\to Y\to Z\dashrightarrow$ for such a conflation.
The extension bifunctor and projective objects are
\begin{equation}\label{eq:extriangulated}
 \mathbb E_{\H_A^m}(Z,X)=\Hom_{\D_{\fd}(A)}(Z,X[1]),
 \qquad\mathcal P(\H_A^m)=\add A;
\end{equation}
see \cite[\S1 and Corollary 2.3]{MP25}.
For $\mathcal S\subseteq\H_A^m$, define
$\Filt_{\H_A^m}(\mathcal S)$ using finite sequences of conflations
\begin{equation}\label{eq:filtration}
 M_{a-1}\longrightarrow M_a\longrightarrow F_a\dashrightarrow,
 \qquad M_0=0,\quad M_\ell=M,\quad F_a\in\mathcal S
 \quad(1\leq a\leq\ell).
\end{equation}
All intermediate objects $M_a$ belong to $\H_A^m$.
The empty filtration is permitted. Closure under direct summands is
not part of the definition of $\Filt_{\H_A^m}(\mathcal S)$.

\begin{definition}[$m$-qh dga]\label{def:dq}
An algebra $A$ satisfying the standing assumptions is an
\defn{$m$-qh dga with respect to $\sigma$} if there are objects
$\std_v\in\H_A^m$, called \defn{standard objects}, satisfying:
\begin{enumerate}[label=\textup{(dq\arabic*)},ref=dq\arabic*]
\item\label{ax:dq1} For $1\leq v,w\leq n$ and $q\in\ZZ$,
\begin{equation}\label{eq:dq1}
 \begin{gathered}
 H^q(\std_v)e_w=0\quad\text{if }v<_\sigma w,\\
 H^q(\std_v)e_v\cong
 \begin{cases}\Bbbk,&q=0,\\0,&q\neq0.
 \end{cases}
 \end{gathered}
\end{equation}
\item\label{ax:dq2} For every vertex $v$, there is a conflation
\begin{equation}\label{eq:presentation}
 K_v\longrightarrow P_v\xrightarrow{p_v}\std_v\dashrightarrow
 \qquad\text{in }\H_A^m,
\end{equation}
with
\begin{equation}\label{eq:kernel}
 \begin{aligned}
 K_v&\in\Filt_{\H_A^m}(\mathcal S_{>v}^{m}),\\
 \mathcal S_{>v}^{m}
 &=\{\std_w[s]:v<_\sigma w,\ 0\leq s<m,
                         \ \std_w[s]\in\H_A^m\}.
 \end{aligned}
\end{equation}
\end{enumerate}
\end{definition}

We call the parameter $m$ in Definition~\ref{def:dq} the
\defn{width} of the structure. An integer $\ell\geq1$ is an
\defn{admissible width} for $(A,\sigma)$ if the standing assumptions
and Definition~\ref{def:dq} hold with $m=\ell$.

Condition~\ref{ax:dq1} says that $H^0(\std_v)$ has only composition
factors $S_w$ with $w\leq_\sigma v$, with $S_v$ occurring once,
whereas $H^q(\std_v)$ for $q<0$ has only composition factors $S_w$
with $w<_\sigma v$. Condition~\ref{ax:dq2} requires every shifted
factor, as well as every intermediate object, to lie in $\H_A^m$.
The object $K_v$ is the first term of the conflation
\eqref{eq:presentation}; no assertion that $K_v$ is an ordinary
module kernel is intended. For the greatest vertex $\sigma(n)$,
\begin{equation}\label{eq:maximal-standard}
 K_{\sigma(n)}=0,\qquad p_{\sigma(n)}:
 P_{\sigma(n)}\xrightarrow{\sim}\std_{\sigma(n)}.
\end{equation}
Brickness, exceptionality, fullness, and smoothness are not additional
requirements in Definition~\ref{def:dq}.

\section{Perfection, fullness, and smoothness}\label{sec:smoothness}

\begin{proposition}[The consequence of the filtration axiom alone]
\label{prop:perfection}
For objects $\std_v\in\H_A^m$ satisfying condition~\ref{ax:dq2},
\begin{equation}\label{eq:perfect-generation}
 \std_v\in\per A\quad(1\leq v\leq n),\qquad
 \thick_{\D_{\fd}(A)}(\std_1,\ldots,\std_n)=\per A.
\end{equation}
For every vertex $v$ one also has
\begin{equation}\label{eq:tail-generation}
 \thick_{\D_{\fd}(A)}\{P_w:v<_\sigma w\}
 =\thick_{\D_{\fd}(A)}\{\std_w:v<_\sigma w\}.
\end{equation}
\end{proposition}
\begin{proof}
Start with \eqref{eq:maximal-standard} and descend in $\leq_\sigma$.
If $\std_w$ is perfect for all $w>_\sigma v$, the filtration
\eqref{eq:kernel} makes $K_v$ perfect; the triangle
\eqref{eq:presentation} then makes $\std_v$ perfect.
Conversely, \eqref{eq:presentation} and \eqref{eq:kernel} express
every $P_v$ in the thick subcategory generated by the standards.
Restricting both arguments to vertices $w>_\sigma v$ proves
\eqref{eq:tail-generation}.
\end{proof}

\begin{theorem}[Smoothness is a consequence]\label{thm:smoothness}
Conditions~\ref{ax:dq1} and \ref{ax:dq2} imply
\begin{equation}\label{eq:smoothness}
 \begin{gathered}
 S_v\in\per A\quad(1\leq v\leq n),\\
 \D_{\fd}(A)=\per A
 =\thick_{\D_{\fd}(A)}(\std_1,\ldots,\std_n),\\
 A\in\per(A^{\op}\otimes_{\Bbbk}A).
 \end{gathered}
\end{equation}
The condition $A\in\per(A^{\op}\otimes_{\Bbbk}A)$ defines
\defn{homological smoothness}.
\end{theorem}
\begin{proof}
\noindent\textit{Step 1: perfect standards.}
Proposition~\ref{prop:perfection} gives $\std_v\in\per A$.

\noindent\textit{Step 2: perfect simples.}
Since $H^1(K_v)=0$, the morphism
$H^0(p_v):e_vB\twoheadrightarrow H^0(\std_v)$ is surjective.
Condition~\ref{ax:dq1} gives $[H^0(\std_v):S_v]=1$.
Thus $H^0(\std_v)$ has simple top $S_v$. Choose a top quotient
$q_v:H^0(\std_v)\twoheadrightarrow S_v$, and let
$a_v:\std_v\to S_v$ be the composite of
$\std_v\to H^0(\std_v)$ with $q_v$. Take the triangle
\begin{equation}\label{eq:simple-triangle}
 T_v\longrightarrow\std_v\xrightarrow{a_v}S_v
 \longrightarrow T_v[1].
\end{equation}
The cohomology sequence of \eqref{eq:simple-triangle} gives
\begin{equation}\label{eq:simple-fibre}
 H^q(T_v)\cong
 \begin{cases}
 H^q(\std_v),&q<0,\\
 \Ker q_v,&q=0,\\
 0,&q>0.
 \end{cases}
\end{equation}
Every composition factor of every $H^q(T_v)$ is $S_w$ with
$w<_\sigma v$, by condition~\ref{ax:dq1}. Finite truncation
triangles and composition series give
\begin{equation}\label{eq:simple-induction}
 T_v\in\thick_{\D_{\fd}(A)}\{S_w:w<_\sigma v\}.
\end{equation}
For the least vertex $\sigma(1)$, $T_{\sigma(1)}=0$ and
$S_{\sigma(1)}\cong\std_{\sigma(1)}$ is perfect. Ascending
induction using \eqref{eq:simple-triangle} makes every $S_v$ perfect.

\noindent\textit{Step 3: fullness and smoothness.}
Every object of $\D_{\fd}(A)$ is obtained from the $S_v$ by finite
truncation triangles and composition series. Consequently
$\D_{\fd}(A)\subseteq\per A$. Properness gives the reverse
inclusion. Finally, $B/\rad B\cong\Bbbk^n$ is separable over
$\Bbbk$ and perfect over $A$. The criterion
\cite[Theorem 5.8]{RS22} proves
$A\in\per(A^{\op}\otimes_{\Bbbk}A)$.
\end{proof}

\section{Exceptionality and finite Hom tests}\label{sec:exceptionality}

An object $E\in\D_{\fd}(A)$ is \defn{exceptional} if
\begin{equation}\label{eq:exceptional}
 \Hom_{\D_{\fd}(A)}(E,E[t])\cong
 \begin{cases}\Bbbk,&t=0,\\0,&t\neq0\end{cases}
 \qquad(t\in\ZZ).
\end{equation}
An ordered tuple $(E_1,\ldots,E_n)$ is an
\defn{exceptional sequence} if every $E_a$ is exceptional and
\begin{equation}\label{eq:exceptional-sequence}
 \Hom_{\D_{\fd}(A)}(E_a,E_b[t])=0
 \qquad(1\leq b<a\leq n,\ t\in\ZZ).
\end{equation}
The sequence is \defn{full} if its terms thickly generate
$\D_{\fd}(A)$. We use the ordering
\begin{equation}\label{eq:standard-sequence}
 \boldsymbol\std^\sigma=(\std_{\sigma(1)},\ldots,\std_{\sigma(n)}).
\end{equation}
Thus standards indexed by larger vertices in $\leq_\sigma$ occur
later in \eqref{eq:standard-sequence}, and backward orthogonality means
vanishing from $\std_w$ to $\std_v[t]$ for $w>_\sigma v$.

\begin{lemma}[Orthogonality from composition factors]
\label{prop:orthogonality}
Assume condition~\ref{ax:dq2} and
\begin{equation}\label{eq:off-diagonal-support}
 H^t(\std_v)e_w=0\qquad(v<_\sigma w,\ t\in\ZZ).
\end{equation}
Then
\begin{align}
 \Hom_{\D_{\fd}(A)}(\std_w,\std_v[t])&=0
 &&(v<_\sigma w,\ t\in\ZZ),\label{eq:backward}\\
 p_v^*:\Hom_{\D_{\fd}(A)}(\std_v,\std_v[t])
 &\xrightarrow{\sim}\Hom_{\D_{\fd}(A)}(P_v,\std_v[t])
 \cong H^t(\std_v)e_v
 &&(1\leq v\leq n,\ t\in\ZZ).
 \label{eq:diagonal}
\end{align}
\end{lemma}
\begin{proof}
Fix $v$ and consider the vertices $w>_\sigma v$ in decreasing
order with respect to $\leq_\sigma$.
Equation~\eqref{eq:projective-cohomology} and
\eqref{eq:off-diagonal-support} give
$\Hom_{\D_{\fd}(A)}(P_w,\std_v[t])=0$ for every $t$.
Once every standard $\std_u$ with $u>_\sigma w$ is orthogonal to
$\std_v$ in every degree, the filtration of $K_w$ in
\eqref{eq:kernel} gives the same vanishing for $K_w$.
Apply $\Hom_{\D_{\fd}(A)}(-,\std_v[t])$ to the presentation
\eqref{eq:presentation} for $w$ to obtain \eqref{eq:backward}.
The filtration of $K_v$ now implies
$\Hom_{\D_{\fd}(A)}(K_v,\std_v[t])=0$ for every $t$.
The presentation for $v$ gives \eqref{eq:diagonal}.
\end{proof}

\begin{theorem}[Equivalent formulations under the filtration axiom]
\label{thm:equivalence}
For objects $\std_v\in\H_A^m$ satisfying condition~\ref{ax:dq2},
the following are equivalent:
\begin{enumerate}[label=\textup{(\roman*)}]
\item condition~\ref{ax:dq1};
\item $\boldsymbol\std^\sigma$ is an exceptional sequence;
\item $\boldsymbol\std^\sigma$ is full exceptional in
      $\D_{\fd}(A)=\per A$.
\end{enumerate}
In particular, $\End_{\H_A^m}(\std_v)\cong\Bbbk$ is a consequence
of conditions~\ref{ax:dq1} and \ref{ax:dq2}.
\end{theorem}
\begin{proof}
We prove $(\mathrm{i})\Rightarrow(\mathrm{ii})\Rightarrow
(\mathrm{i})$, then $(\mathrm{i})\Rightarrow(\mathrm{iii})$;
$(\mathrm{iii})\Rightarrow(\mathrm{ii})$ follows from the definition.

\noindent\underline{$(\mathrm{i})\Rightarrow(\mathrm{ii})$.}
Lemma~\ref{prop:orthogonality} gives \eqref{eq:backward}.
Equations \eqref{eq:diagonal} and \eqref{eq:dq1} give
\eqref{eq:exceptional} for every $\std_v$.

\noindent\underline{$(\mathrm{ii})\Rightarrow(\mathrm{i})$.}
For $w>_\sigma v$, equation~\eqref{eq:tail-generation} places $P_w$
in the thick subcategory generated by $\std_u$ with $u>_\sigma v$.
Equation~\eqref{eq:exceptional-sequence} gives
$\Hom_{\D_{\fd}(A)}(P_w,\std_v[t])=0$ for all $t$.
Thus \eqref{eq:off-diagonal-support} holds.
The isomorphism \eqref{eq:diagonal} and exceptionality of
$\std_v$ give $H^t(\std_v)e_v\cong\Bbbk$ for $t=0$ and zero for
$t\neq0$, completing \eqref{eq:dq1}.

\noindent\underline{$(\mathrm{i})\Rightarrow(\mathrm{iii})$.}
Theorem~\ref{thm:smoothness} gives fullness and
$\D_{\fd}(A)=\per A$; the implication
$(\mathrm{i})\Rightarrow(\mathrm{ii})$ gives exceptionality.
\end{proof}

\begin{proposition}[Tests inside the extended module category]
\label{prop:finite-tests}
For every vertex $w$ and $0\leq s<m$, define
\begin{equation}\label{eq:truncated-projectives}
 U_{w,s}=\tau^{\geq1-m}(P_w[s])\in\H_A^m.
\end{equation}
There are natural isomorphisms
\begin{equation}\label{eq:intrinsic-tests}
 \Hom_{\H_A^m}(U_{w,s},\std_v)
 \cong\Hom_{\D_{\fd}(A)}(P_w,\std_v[-s])
 \cong H^{-s}(\std_v)e_w.
\end{equation}
Consequently condition~\ref{ax:dq1} is equivalent to
\begin{equation}\label{eq:finite-tests}
 \begin{aligned}
 \Hom_{\H_A^m}(U_{w,s},\std_v)&=0
 &&(v<_\sigma w,\ 0\leq s<m),\\
 \Hom_{\H_A^m}(U_{v,s},\std_v)&=0
 &&(1\leq s<m),\\
 \dim_{\Bbbk}\Hom_{\H_A^m}(P_v,\std_v)&=1
 &&(1\leq v\leq n).
 \end{aligned}
\end{equation}
Under condition~\ref{ax:dq2} and the first line of
\eqref{eq:finite-tests}, the last line can equivalently be replaced
by the \defn{brick condition}
$\End_{\H_A^m}(\std_v)\cong\Bbbk$ for every $v$.
\end{proposition}
\begin{proof}
Put $L_{w,s}=\tau^{\leq-m}(P_w[s])$. In the triangle
\[
 L_{w,s}\longrightarrow P_w[s]\longrightarrow U_{w,s}
 \longrightarrow L_{w,s}[1],
\]
both $L_{w,s}$ and $L_{w,s}[1]$ lie in
$\D_{\fd}(A)^{\leq-m}$, whereas $\std_v$ lies in
$\D_{\fd}(A)^{\geq1-m}$. Applying
$\Hom_{\D_{\fd}(A)}(-,\std_v)$ proves
\eqref{eq:intrinsic-tests}. Since all other cohomological degrees
of $\std_v$ vanish, \eqref{eq:finite-tests} is equivalent to
\eqref{eq:dq1}. Under condition~\ref{ax:dq2},
\eqref{eq:diagonal} identifies the last vector space in
\eqref{eq:finite-tests} with $\End_{\H_A^m}(\std_v)$.
\end{proof}

\begin{corollary}[Classical specialisation]\label{cor:classical}
For $m=1$, Definition~\ref{def:dq} is equivalent to classical
quasi-heredity of $B=H^0(A)$ with respect to $\leq_\sigma$.
\end{corollary}
\begin{proof}
For $m=1$, the dga $A$ is quasi-isomorphic to $B$, and
$\H_A^1\simeq\mod B$. We prove both implications.

\noindent\underline{Conditions \ref{ax:dq1} and \ref{ax:dq2}
$\Rightarrow$ classical quasi-heredity.}
Only the shift $s=0$ is allowed in \eqref{eq:kernel}, and
\eqref{eq:presentation} becomes
$0\to K_v\to e_vB\to\std_v\to0$.
Condition~\ref{ax:dq1} gives composition factors $S_w$ with
$w\leq_\sigma v$, with $S_v$ occurring once. These are the classical projective
presentation axioms \cite[\S2.1]{BB26}.

\noindent\underline{Classical quasi-heredity $\Rightarrow$
conditions \ref{ax:dq1} and \ref{ax:dq2}.}
The classical standard modules have the required composition factors,
and the projective kernel $K_v$ is filtered by $\std_w$ with
$w>_\sigma v$.
Thus \eqref{eq:dq1} and \eqref{eq:kernel} hold with $s=0$.
\end{proof}

\section{Canonical standards and costandards}\label{adv:canonical-section}

Assume conditions~\ref{ax:dq1} and \ref{ax:dq2} throughout
Sections~\ref{adv:canonical-section}--\ref{adv:opposite-section}.

\begin{proposition}[Standards determined by the ordered projectives]
\label{adv:canonical}
For a vertex $v$, set
\begin{equation}\label{adv:tail}
 \mathcal Q_v=\thick_{\D_{\fd}(A)}\{P_w:v<_\sigma w\},
 \qquad e_{>v}^\sigma=\sum_{v<_\sigma w}e_w.
\end{equation}
Define $\mathcal Q_v^\perp$ to consist of the objects $X\in\D_{\fd}(A)$
satisfying $\Hom_{\D_{\fd}(A)}(Y,X[t])=0$ for every
$Y\in\mathcal Q_v$ and $t\in\ZZ$. Then
\begin{equation}\label{adv:localisation}
 \begin{gathered}
 K_v\in\mathcal Q_v,\qquad \std_v\in\mathcal Q_v^\perp,\\
 p_v^*:\Hom_{\D_{\fd}(A)}(\std_v,X[t])
 \xrightarrow{\sim}\Hom_{\D_{\fd}(A)}(P_v,X[t])
 \quad(X\in\mathcal Q_v^\perp,\ t\in\ZZ).
 \end{gathered}
\end{equation}
The pair $(\std_v,p_v)$ is determined by $(A,\sigma)$ up to a unique
isomorphism compatible with $p_v$, and $K_v$ is determined up to
isomorphism. Moreover,
\begin{equation}\label{adv:ordinary-standard}
 H^0(\std_v)\cong e_vB/e_vBe_{>v}^\sigma B,
 \qquad B=H^0(A).
\end{equation}
The morphism $H^0(p_v):e_vB\to H^0(\std_v)$ is a projective cover
with simple top $S_v$.
\end{proposition}
\begin{proof}
Proposition~\ref{prop:perfection} identifies $\mathcal Q_v$ with
$\thick_{\D_{\fd}(A)}\{\std_w:v<_\sigma w\}$.
The filtration of $K_v$ therefore gives $K_v\in\mathcal Q_v$.
Condition~\ref{ax:dq1} gives
$\Hom_{\D_{\fd}(A)}(P_w,\std_v[t])=H^t(\std_v)e_w=0$ for
$w>_\sigma v$, so $\std_v\in\mathcal Q_v^\perp$.
Applying $\Hom_{\D_{\fd}(A)}(-,X[t])$ to the presentation triangle
\eqref{eq:presentation} proves \eqref{adv:localisation}.
The universal property \eqref{adv:localisation} proves uniqueness of
$(\std_v,p_v)$; cones then give uniqueness of $K_v$ up to isomorphism.

Since $H^1(K_v)=0$, the morphism $H^0(p_v)$ is surjective.
Put $D_v^0=e_vB/e_vBe_{>v}^\sigma B$. The support condition in
\ref{ax:dq1} makes $H^0(p_v)$ factor through a surjection
$f_v:D_v^0\to H^0(\std_v)$.
The module $D_v^0$ belongs to $\mathcal Q_v^\perp$, so its quotient
morphism $P_v\to D_v^0$ factors uniquely through $p_v$.
Taking degree-zero cohomology gives
$g_v:H^0(\std_v)\to D_v^0$.
Precomposition with the surjections from $e_vB$ shows
$g_vf_v=1_{D_v^0}$ and $f_vg_v=1_{H^0(\std_v)}$.
This proves \eqref{adv:ordinary-standard}.
Condition~\ref{ax:dq1} gives $H^0(\std_v)e_v\cong\Bbbk$, so
$H^0(\std_v)$ is a non-zero quotient of the indecomposable projective
$e_vB$. The top of $H^0(\std_v)$ is $S_v$, and
$H^0(p_v):e_vB\to H^0(\std_v)$ is a projective cover.
\end{proof}

\begin{definition}[Costandard sequence]\label{adv:dual-definition}
The \defn{left-dual sequence} of
$\boldsymbol\std^\sigma$
is the exceptional sequence
$(\costd_{\sigma(n)},\ldots,\costd_{\sigma(1)})$ satisfying
\begin{equation}\label{adv:pairing}
 \Hom_{\D_{\fd}(A)}(\std_u,\costd_v[t])\cong
 \begin{cases}
 \Bbbk,&u=v\text{ and }t=0,\\
 0,&\text{otherwise}
 \end{cases}
 \qquad(1\leq u,v\leq n,\ t\in\ZZ).
\end{equation}
The objects $\costd_v$ are the \defn{costandards}.
The ordering and pairing in Definition~\ref{adv:dual-definition}
specify the convention for left-dual sequences; compare
\cite[\S2.3]{BB26}.
\end{definition}

\begin{proposition}[Existence, support, and amplitude of costandards]
\label{adv:costandards}
The left-dual sequence in Definition~\ref{adv:dual-definition} exists
and is unique up to isomorphism of its terms. It is full in
$\D_{\fd}(A)=\per A$, and
\begin{equation}\label{adv:costandard-window}
 \costd_v\in\H_A^m[1-m]
 =\D_{\fd}(A)^{[0,m-1]}.
\end{equation}
The vertex components of each $\costd_v$ satisfy
\begin{equation}\label{adv:costandard-support}
 \begin{gathered}
 H^t(\costd_v)e_w=0\quad(v<_\sigma w,\ t\in\ZZ),\\
 H^t(\costd_v)e_v\cong
 \begin{cases}\Bbbk,&t=0,\\0,&t\neq0.\end{cases}
 \end{gathered}
\end{equation}
\end{proposition}
\begin{proof}
\noindent\textit{Step 1: construction and uniqueness.}
For $v=\sigma(a)$, start with $Y_a=\std_v$. For
$j=a-1,a-2,\ldots,1$, form evaluation triangles
\begin{equation}\label{adv:mutation}
 \begin{gathered}
 V_j=\bigoplus_{t\in\ZZ}
 \Hom_{\D_{\fd}(A)}(\std_{\sigma(j)},Y_{j+1}[t])[-t],\\
 V_j\otimes_{\Bbbk}\std_{\sigma(j)}\longrightarrow
 Y_{j+1}\longrightarrow Y_j\longrightarrow
 (V_j\otimes_{\Bbbk}\std_{\sigma(j)})[1].
 \end{gathered}
\end{equation}
The complexes $V_j$ have zero differential and finite total dimension:
all standards are perfect, and every $Y_{j+1}$ has finite-dimensional
total cohomology. Put $\costd_v=Y_1$; for $a=1$, put
$\costd_v=\std_v$.
Exceptionality annihilates morphisms from $\std_{\sigma(j)}$ at step
$j$, and backward orthogonality preserves the vanishings already
obtained. Thus $\costd_v$ lies in
$\thick_{\D_{\fd}(A)}\{\std_u:u\leq_\sigma v\}$ and is right
orthogonal to every $\std_u$ with $u<_\sigma v$.
The same triangles give \eqref{adv:pairing}.
Combining the mutation triangles gives a triangle
$M_v\to\std_v\to\costd_v\to M_v[1]$, with
$M_v\in\thick_{\D_{\fd}(A)}\{\std_u:u<_\sigma v\}$.
Consequently
\[
 \Hom_{\D_{\fd}(A)}(\costd_v,\costd_v[t])
 \cong\Hom_{\D_{\fd}(A)}(\std_v,\costd_v[t]),
\]
so $\costd_v$ is exceptional.
For $u<_\sigma v$, the object $\costd_u$ lies in the subcategory
generated by $\std_z$ with $z<_\sigma v$, giving
$\Hom_{\D_{\fd}(A)}(\costd_u,\costd_v[t])=0$.
Induction in $a$ using the triangles for $M_v$ shows that the
costandards generate every standard, hence gives fullness of the
reversed sequence.

For uniqueness, an object $N_v\in\D_{\fd}(A)$ satisfying \eqref{adv:pairing}
belongs to $\thick_{\D_{\fd}(A)}\{\std_u:u\leq_\sigma v\}$:
evaluation against $\boldsymbol\std^\sigma$ has no factors belonging
to $\thick_{\D_{\fd}(A)}(\std_w)$ with $w>_\sigma v$.
A generator of $\Hom_{\D_{\fd}(A)}(\std_v,N_v)\cong\Bbbk$ gives a
triangle $M_v\to\std_v\to N_v\to M_v[1]$, where
$M_v\in\thick_{\D_{\fd}(A)}\{\std_u:u<_\sigma v\}$.
Since $N_v$ is right orthogonal to that subcategory, its localisation
property identifies $N_v$ with the object constructed in
\eqref{adv:mutation}.

\noindent\textit{Step 2: the cohomological window.}
For any allowed kernel factor, \eqref{adv:pairing} gives
\begin{equation}\label{adv:factor-pairing}
 \Hom_{\D_{\fd}(A)}(\std_w[s],\costd_u[t])\cong
 \begin{cases}\Bbbk,&w=u\text{ and }t=s,\\0,&\text{otherwise}.
 \end{cases}
\end{equation}
Since $0\leq s<m$, the filtration of $K_v$ implies
$\Hom_{\D_{\fd}(A)}(K_v,\costd_u[t])=0$ for $t\notin[0,m-1]$.
Apply $\Hom_{\D_{\fd}(A)}(-,\costd_u[t])$ to
\eqref{eq:presentation}; equation~\eqref{adv:pairing} then gives
$H^t(\costd_u)e_v=0$ for every vertex $v$ and $t\notin[0,m-1]$.
This proves \eqref{adv:costandard-window}.

\noindent\textit{Step 3: vertex support.}
The factors in a standard filtration of $P_w$, supplied by
condition~\ref{ax:dq2}, are shifts of $\std_z$ with $z\geq_\sigma w$.
If $w>_\sigma v$, all their morphisms to shifts of
$\costd_v$ vanish by \eqref{adv:factor-pairing}.
If $w=v$, only the quotient $\std_v$ contributes, in degree zero.
Since $H^t(\costd_v)e_w=
\Hom_{\D_{\fd}(A)}(P_w,\costd_v[t])$, the calculations for
$w>_\sigma v$ and $w=v$ prove
\eqref{adv:costandard-support}.
\end{proof}

\section{Filtrations and the least possible width}\label{adv:filtration-section}

For $\ell\geq1$, define
\begin{equation}\label{adv:allowed}
 \begin{gathered}
 I_w^\ell=\{s\in\ZZ:0\leq s<\ell,\ \std_w[s]\in\H_A^\ell\},\\
 \mathcal S_{>v}^\ell
 =\{\std_w[s]:v<_\sigma w,\ s\in I_w^\ell\}.
 \end{gathered}
\end{equation}
Here $\H_A^\ell=\D_{\fd}(A)^{[1-\ell,0]}$, whether or not $A$ has
cohomology concentrated in that interval.

\begin{proposition}[Recognition of shifted standard filtrations]
\label{adv:criterion}
For $X\in\mathcal Q_v$, the following conditions are equivalent:
\begin{enumerate}[label=\textup{(\arabic*)}]
\item\label{adv:criterion-filtration}
$X\in\Filt_{\H_A^\ell}(\mathcal S_{>v}^\ell)$;
\item\label{adv:criterion-vanishing}
$\Hom_{\D_{\fd}(A)}(X,\costd_w[t])=0$ for every
$w>_\sigma v$ and $t\notin I_w^\ell$.
\end{enumerate}
If condition~\ref{adv:criterion-vanishing} holds, put
$a=\sigma^{-1}(v)$. The object $X$ admits a filtration with successive
factors $G_{\sigma(n)}(X),\ldots,G_{\sigma(a+1)}(X)$, where
\begin{equation}\label{adv:standard-factors}
 G_w(X)\cong\bigoplus_{s\in I_w^\ell}
 \std_w[s]^{\oplus\dim_{\Bbbk}
 \Hom_{\D_{\fd}(A)}(X,\costd_w[s])}.
\end{equation}
In particular, $\Filt_{\H_A^\ell}(\mathcal S_{>v}^\ell)$ is closed
under direct summands.
\end{proposition}
\begin{proof}
\noindent\underline{$(1)\Rightarrow(2)$.}
Every allowed factor satisfies the required vanishings by
\eqref{adv:factor-pairing}; the long exact sequences of the filtration
preserve the vanishings in condition~\ref{adv:criterion-vanishing}.

\noindent\underline{$(2)\Rightarrow(1)$.}
Put $a=\sigma^{-1}(v)$. The full exceptional sequence in
$\mathcal Q_v$ is
$(\std_{\sigma(a+1)},\ldots,\std_{\sigma(n)})$.
Successive evaluation against
$\std_{\sigma(n)},\ldots,\std_{\sigma(a+1)}$ gives a tower with
one factor $G_w(X)\in\thick_{\D_{\fd}(A)}(\std_w)$ for each
$w>_\sigma v$.
Equivalently, the octahedral axiom gives a finite sequence of triangles
starting at zero and ending at $X$, with successive factors
$G_{\sigma(n)}(X),\ldots,G_{\sigma(a+1)}(X)$.
Exceptionality expresses $G_w(X)$ as
a finite direct sum of shifts of $\std_w$.
For a fixed $w$, applying
$\Hom_{\D_{\fd}(A)}(-,\costd_w[t])$ to the tower annihilates every
factor $G_u(X)$ with $u\neq w$. Consequently the multiplicity
of $\std_w[t]$ in $G_w(X)$ is
$\dim_{\Bbbk}\Hom_{\D_{\fd}(A)}(X,\costd_w[t])$.
Condition~\ref{adv:criterion-vanishing} gives \eqref{adv:standard-factors}.
Every factor in \eqref{adv:standard-factors} belongs to $\H_A^\ell$;
extension closure puts every intermediate object in $\H_A^\ell$.
Refining the finite direct sums gives
condition~\ref{adv:criterion-filtration}.
When $v=\sigma(n)$, the category $\mathcal Q_v$ is zero, and the
empty filtration proves the assertion.

Both $\mathcal Q_v$ and the vanishings in
condition~\ref{adv:criterion-vanishing} are closed
under direct summands, proving closure of
$\Filt_{\H_A^\ell}(\mathcal S_{>v}^\ell)$ under direct summands.
\end{proof}

For vertices $v,w$ and $s\in\ZZ$, set
\begin{equation}\label{adv:multiplicities}
 b_{vw,s}=\dim_{\Bbbk}\Hom_{\D_{\fd}(A)}(P_v,\costd_w[s])
         =\dim_{\Bbbk}H^s(\costd_w)e_v.
\end{equation}
For $v<_\sigma w$, applying
$\Hom_{\D_{\fd}(A)}(-,\costd_w[s])$ to
\eqref{eq:presentation} gives
\begin{equation}\label{adv:kernel-pairing}
 \Hom_{\D_{\fd}(A)}(K_v,\costd_w[s])
 \cong\Hom_{\D_{\fd}(A)}(P_v,\costd_w[s]).
\end{equation}
Thus the ordered filtration of $K_v$ in
Proposition~\ref{adv:criterion} contains exactly $b_{vw,s}$ copies of
$\std_w[s]$ in the factor $G_w(K_v)$.

\begin{example}[Unsigned filtration multiplicities are not invariant]
\label{adv:zero-filtration}
Suppose $w>_\sigma v$ and both $\std_w[s]$ and $\std_w[s+1]$ belong
to $\mathcal S_{>v}^m$. The triangles
\begin{equation}\label{adv:cancellation}
 0\longrightarrow\std_w[s]\xrightarrow{1}\std_w[s]\longrightarrow0,
 \qquad
 \std_w[s]\longrightarrow0\longrightarrow\std_w[s+1]
 \xrightarrow{1}\std_w[s+1]
\end{equation}
define a two-factor filtration of the zero object, with factors
$\std_w[s]$ and $\std_w[s+1]$.
For a concrete instance, take $A=\Bbbk\times\Bbbk$, $m=2$,
$1<_\sigma2$, $\std_i=P_i=S_i$, $w=2$, and $s=0$.
Therefore arbitrary filtrations need not have invariant lengths or
invariant numbers of shifted factors, even for a semisimple algebra.
The multiplicities in the ordered factors \eqref{adv:standard-factors}
are nevertheless invariant.
\end{example}

\begin{corollary}[Least width]\label{adv:width}
Suppose $(A,\sigma)$ admits conditions~\ref{ax:dq1} and
\ref{ax:dq2} for at least one width. For its canonical standards and
costandards, put
\begin{equation}\label{adv:depths}
 \begin{aligned}
 \alpha_A&=-\min\{q:H^q(A)\neq0\},\\
 \alpha_w&=-\min\{q:H^q(\std_w)\neq0\},
 &\beta_w&=\max\{q:H^q(\costd_w)\neq0\}.
 \end{aligned}
\end{equation}
These integers are non-negative. The least admissible width is
\begin{equation}\label{adv:least-width}
 m_{\min}(A,\sigma)=
 1+\max\left\{\alpha_A,\
             \max_{1\leq w\leq n}(\alpha_w+\beta_w)\right\}.
\end{equation}
Every width $m\geq m_{\min}(A,\sigma)$ is admissible.
\end{corollary}
\begin{proof}
Propositions~\ref{adv:canonical} and \ref{adv:costandards} make the
objects in \eqref{adv:depths} independent of the chosen width.
$H^0(A)$ is non-zero. Condition~\ref{ax:dq1} and
\eqref{adv:costandard-support} give
$H^0(\std_w)e_w\cong H^0(\costd_w)e_w\cong\Bbbk$ for every
vertex $w$. Hence all the integers
in \eqref{adv:depths} are non-negative.
For $s\geq0$, the cohomological shift convention gives
\begin{equation}\label{adv:shift-window}
 \std_w[s]\in\H_A^\ell
 \quad\Longleftrightarrow\quad
 0\leq s\leq\ell-1-\alpha_w.
\end{equation}
If $\beta_w>0$, equation~\eqref{adv:costandard-support} gives a
vertex $v<_\sigma w$ with $b_{vw,\beta_w}>0$.
Necessity in Proposition~\ref{adv:criterion}, applied to $K_v$,
then forces $\alpha_w+\beta_w\leq\ell-1$ for every admissible width
$\ell$. If $\beta_w=0$, the same inequality follows from
$\std_w\in\H_A^\ell$. Also $\alpha_A\leq\ell-1$ is necessary.

Conversely, assume $\alpha_A\leq\ell-1$ and
$\alpha_w+\beta_w\leq\ell-1$ for every $w$. The objects $P_v$ and
$\std_v$ lie in $\H_A^\ell$. Every shift occurring in the ordered
factors $G_w(K_v)$ satisfies $s\leq\beta_w$; equation
\eqref{adv:shift-window} therefore puts each factor in $\H_A^\ell$.
Proposition~\ref{adv:criterion} gives the required kernel filtration,
including $K_v\in\H_A^\ell$.
The presentation triangle is a conflation in $\H_A^\ell$, and
condition~\ref{ax:dq1} is unchanged. This proves
\eqref{adv:least-width} and the assertion about larger widths.
\end{proof}

\begin{corollary}[Standards and costandards concentrated in degree zero]
\label{adv:both-hearts}
The following are equivalent:
\begin{enumerate}[label=\textup{(\roman*)}]
\item $\std_v,\costd_v\in\mod B$ for every vertex $v$;
\item $m_{\min}(A,\sigma)=1$.
\end{enumerate}
In either case, $A$ is quasi-isomorphic to the ordinary
quasi-hereditary algebra $B$ with the given order.
\end{corollary}
\begin{proof}
\noindent\underline{$(\mathrm{i})\Rightarrow(\mathrm{ii})$.}
Degree-zero concentration of every $\costd_w$ gives $b_{vw,s}=0$
for $s\neq0$ in \eqref{adv:multiplicities}.
Proposition~\ref{adv:criterion} therefore gives unshifted standard
filtrations of the kernels $K_v$. Since every $\std_w$ is a module,
every $K_v$ and then every $P_v$ belongs to $\mod B$.
Thus $A\simeq B$ and all three integers in
\eqref{adv:depths} are zero. Equation~\eqref{adv:least-width}
gives $m_{\min}(A,\sigma)=1$.

\noindent\underline{$(\mathrm{ii})\Rightarrow(\mathrm{i})$.}
Apply Definition~\ref{def:dq} and
Proposition~\ref{adv:costandards} at width $1$.
Corollary~\ref{cor:classical} gives classical quasi-heredity.
\end{proof}

\section{Grothendieck groups, injectives, and opposite algebras}
\label{adv:opposite-section}

The \defn{Grothendieck group} $K_0(\D_{\fd}(A))$ is generated by
symbols $[X]$, with relations $[Y]=[X]+[Z]$ for triangles
$X\to Y\to Z\to X[1]$.
Since $\D_{\fd}(A)=\per A$ and $A$ is proper, the \defn{Euler pairing}
\begin{equation}\label{adv:euler}
 \chi_A(X,Y)=\sum_{t\in\ZZ}(-1)^t
 \dim_{\Bbbk}\Hom_{\D_{\fd}(A)}(X,Y[t])
\end{equation}
is a finite sum and induces a bilinear form on
$K_0(\D_{\fd}(A))$.

\begin{proposition}[Bases, multiplicities, and a determinant]
\label{adv:kzero}
The classes of the simples, standards, costandards, and projectives
are four $\ZZ$-bases of $K_0(\D_{\fd}(A))\cong\ZZ^n$.
The standard and costandard bases satisfy
\begin{equation}\label{adv:dual-bases}
 \chi_A(\std_v,\costd_w)=\delta_{vw}.
\end{equation}
In the projective expansion in the standard basis,
\begin{equation}\label{adv:projective-expansion}
 [P_v]=[\std_v]+\sum_{w>_\sigma v}
           \left(\sum_{s=0}^{m-1}(-1)^s b_{vw,s}\right)[\std_w].
\end{equation}
If $c_{vw,s}$ counts factors $\std_w[s]$ in any filtration of $K_v$,
then
\begin{equation}\label{adv:signed-multiplicity}
 \sum_{s=0}^{m-1}(-1)^s c_{vw,s}
 =\sum_{s=0}^{m-1}(-1)^s b_{vw,s}\qquad(v<_\sigma w).
\end{equation}
Finally, define the \defn{Euler Cartan matrix} by
\begin{equation}\label{adv:cartan}
 C_{vu}=\sum_{q\in\ZZ}(-1)^q[H^q(P_v):S_u]
       =\chi_A(P_u,P_v).
\end{equation}
Then $\det C=1$.
\end{proposition}
\begin{proof}
The bounded standard $t$-structure identifies
$K_0(\D_{\fd}(A))$ with $K_0(\mod B)$, with simple basis $[S_v]$.
Condition~\ref{ax:dq1} gives
$[\std_v]=[S_v]+\sum_{u<_\sigma v}d_{vu}[S_u]$ for integers
$d_{vu}$, so the standard classes are a basis.
Equation~\eqref{adv:costandard-support} proves the analogous assertion
for costandard classes. Equation~\eqref{adv:pairing} gives
\eqref{adv:dual-bases}.
The ordered filtration in Proposition~\ref{adv:criterion} gives
\eqref{adv:projective-expansion}; its transition matrix is upper
unitriangular when rows and columns are listed in $\sigma$-order.
Thus the projective classes form a basis as well.
An arbitrary kernel filtration and the relation
$[\std_w[s]]=(-1)^s[\std_w]$ prove
\eqref{adv:signed-multiplicity} by comparison in the standard basis.
The transition matrix from standards to simples is lower
unitriangular. The product of the projective-to-standard and
standard-to-simple transition matrices is $C$, proving $\det C=1$.
\end{proof}

The integers in \eqref{adv:projective-expansion} and the entries of
\eqref{adv:cartan} can be negative. Proposition~\ref{adv:kzero}
concerns the Euler Cartan matrix of $A$; it does not assert that the
ordinary Cartan matrix of $H^0(A)$ has determinant $1$.

For a complex $V$ of $\Bbbk$-vector spaces, write
$\dual V=\Hom^\bullet_{\Bbbk}(V,\Bbbk)$ for its complex dual.
Vector-space duality exchanges right $A$-modules and right
$A^{\op}$-modules and satisfies
$H^q(\dual X)\cong\dual H^{-q}(X)$.

\begin{proposition}[Injective presentations and opposite algebras]
\label{adv:opposite}
Put $\widehat\costd_v=\costd_v[m-1]$ and
$I_v=\dual(Ae_v)[m-1]$. There are conflations
\begin{equation}\label{adv:injective}
 \begin{gathered}
 \widehat\costd_v\longrightarrow I_v\longrightarrow L_v
 \dashrightarrow\quad\text{in }\H_A^m,\\
 L_v\in\Filt_{\H_A^m}
 \{\widehat\costd_w[-s]:v<_\sigma w,\ 0\leq s<m,\
                    \widehat\costd_w[-s]\in\H_A^m\}.
 \end{gathered}
\end{equation}
The objects $I_v$ are the indecomposable injectives of $\H_A^m$.
The dg algebra $A^{\op}$ satisfies conditions~\ref{ax:dq1} and
\ref{ax:dq2} for the same $m$ and order, with standards
$\dual\costd_v$ and costandards $\dual\std_v$.
In particular,
\begin{equation}\label{adv:opposite-width}
 m_{\min}(A^{\op},\sigma)=m_{\min}(A,\sigma).
\end{equation}
\end{proposition}
\begin{proof}
The shifted vector-space duality
$X\mapsto\dual X[m-1]$ is a contravariant equivalence
$\H_A^m\to\H_{A^{\op}}^m$, taking projectives to injectives.
The projectives of $\H_{A^{\op}}^m$ are the summands of $A$ as a
left $A$-module, so the corresponding right-module injectives are
precisely the objects $I_v$.

Decompose $I_v$ using the full exceptional sequence
$(\costd_{\sigma(n)},\ldots,\costd_{\sigma(1)})$.
The same evaluation-tower argument as in
Proposition~\ref{adv:criterion}, now using
$\Hom_{\D_{\fd}(A)}(\std_w,-)$, gives successive factors in
$\thick_{\D_{\fd}(A)}(\costd_{\sigma(1)}),\ldots,
\thick_{\D_{\fd}(A)}(\costd_{\sigma(n)})$.
The identity
\begin{equation}\label{adv:injective-pairing}
 \Hom_{\D_{\fd}(A)}(\std_w,I_v[t])
 \cong\dual\bigl(H^{1-m-t}(\std_w)e_v\bigr)
\end{equation}
identifies the factor in $\thick_{\D_{\fd}(A)}(\costd_w)$ as
\begin{equation}\label{adv:injective-block}
 \bigoplus_{s\geq0}
 \costd_w[m-1-s]^{\oplus\dim_{\Bbbk}H^{-s}(\std_w)e_v}.
\end{equation}
For $w<_\sigma v$, condition~\ref{ax:dq1} makes the factor
\eqref{adv:injective-block} zero; for $w=v$, the factor is exactly
$\widehat\costd_v$.
Every occurring $s$ satisfies $s\leq\alpha_w$.
Corollary~\ref{adv:width} gives
$\beta_w+s\leq\beta_w+\alpha_w\leq m-1$, so all factors
\eqref{adv:injective-block} belong to $\H_A^m$.
The tower with factors \eqref{adv:injective-block} therefore gives
\eqref{adv:injective}.

Apply $\dual(-)[m-1]$ to the conflation in
\eqref{adv:injective}. The resulting conflation in
$\H_{A^{\op}}^m$ is
\begin{equation}\label{adv:opposite-presentation}
 \dual L_v[m-1]\longrightarrow Ae_v\longrightarrow
 \dual\costd_v\dashrightarrow,
\end{equation}
and its kernel factors are $(\dual\costd_w)[s]$ with
$w>_\sigma v$, $0\leq s<m$, and
$(\dual\costd_w)[s]\in\H_{A^{\op}}^m$.
The support formula \eqref{adv:costandard-support}, after taking
vector-space duals, verifies condition~\ref{ax:dq1} for
$\dual\costd_v$. Thus \eqref{adv:opposite-presentation} proves
condition~\ref{ax:dq2} for $A^{\op}$.
Finally,
\begin{equation}\label{adv:opposite-pairing}
 \Hom_{\D_{\fd}(A^{\op})}(\dual\costd_u,\dual\std_v[t])
 \cong\Hom_{\D_{\fd}(A)}(\std_v,\costd_u[t])
\end{equation}
identifies $\dual\std_v$ as the costandards for $A^{\op}$.
Applying the result to both $A$ and $A^{\op}$ proves
\eqref{adv:opposite-width}.
\end{proof}

\section{Graded path dg algebras and explicit structures}
\label{good:paths}

Let $Q$ be a finite acyclic quiver on $\{1,\ldots,n\}$, with all
arrow degrees at most zero. We concatenate paths from left to right:
$e_ua=a=ae_v$ for an arrow $a:u\to v$.
Let $A=(\Bbbk Q,d_A)$, where $d_A$ is a graded derivation of degree
$1$, $d_A^2=0$, and $d_A(e_v)=0$ for every vertex $v$.
For a path $p$, write $|p|$ for its degree, with $|e_v|=0$, and set
\begin{equation}\label{good:path-width}
 M=1-\min\{|p|:p\text{ is a path in }Q\}.
\end{equation}
The integer $M$ measures the degrees in the underlying graded algebra.
The least integer $m$ satisfying $H^q(A)=0$ for $q\notin[1-m,0]$
can be smaller than $M$ when $d_A\neq0$.

\begin{theorem}[Every total order at the path-degree bound]
\label{good:every-order}
For every $\sigma\in S_n$, put
\begin{equation}\label{good:path-standards}
 e_{>v}^{\sigma}=\sum_{v<_\sigma w}e_w,\qquad
 C_v=A/Ae_{>v}^{\sigma}A,\qquad
 \std_v=e_vC_v,\qquad K_v=e_vAe_{>v}^{\sigma}A.
\end{equation}
The objects in \eqref{good:path-standards}, with their quotient maps
$P_v\twoheadrightarrow\std_v$, satisfy
conditions~\ref{ax:dq1} and~\ref{ax:dq2} for every $m\geq M$.
\end{theorem}
\begin{proof}
\noindent\textit{Step 1: condition~\ref{ax:dq1}.}
The module $\std_v$ has a basis consisting of paths starting at $v$
whose vertices $u$ all satisfy $u\leq_\sigma v$.
Consequently $\std_ve_w=0$ for $v<_\sigma w$.
Acyclicity gives $\std_ve_v=\Bbbk e_v$ in degree zero.
These equalities of complexes imply condition~\ref{ax:dq1}.

\noindent\textit{Step 2: the kernel filtration.}
Fix $v=\sigma(i)$. If $i=n$, then $K_v=0$ and its filtration is
empty. For $i<n$, define dg ideals and dg submodules by
\begin{equation}\label{good:path-tower}
 \begin{gathered}
 J_j=A\left(\sum_{b=j}^n e_{\sigma(b)}\right)A
 \ (i+1\leq j\leq n),\qquad J_{n+1}=0,\\
 F_j=e_vJ_j\quad(i+1\leq j\leq n+1),\\
 0=F_{n+1}\subseteq F_n\subseteq\cdots\subseteq F_{i+1}=K_v.
 \end{gathered}
\end{equation}
For $i<j\leq n$, set $w=\sigma(j)$ and
$L_{vw}=e_vC_we_w$, a finite complex of vector spaces.
Multiplication induces an isomorphism of dg right $A$-modules
\begin{equation}\label{good:path-layer}
 L_{vw}\otimes_{\Bbbk}\std_w\xrightarrow{\sim}F_j/F_{j+1},
 \qquad p\otimes q\longmapsto pq.
\end{equation}
Indeed, a basis path in $F_j/F_{j+1}$ has greatest vertex $w$ with
respect to $\leq_\sigma$. Acyclicity permits a unique cut at its
sole occurrence of $w$. The Leibniz rule makes the multiplication
map in \eqref{good:path-layer} a morphism of complexes.
Splitting the finite complex $L_{vw}$ over $\Bbbk$ yields
\begin{equation}\label{good:path-factors}
 F_j/F_{j+1}\simeq
 \bigoplus_{h=1-M}^{0}
 \std_w[-h]^{\oplus\dim_{\Bbbk}H^h(L_{vw})}
 \quad\text{in }\D_{\fd}(A).
\end{equation}

\noindent\textit{Step 3: admissibility of every filtration factor.}
If $H^h(L_{vw})\neq0$, a prefix path $p:v\to w$ of degree $h$
occurs in $L_{vw}$. For every basis path $q$ in $\std_w$, the
concatenation $pq$ is a path, and $h+|q|\geq1-M$.
Thus $\std_w[-h]\in\H_A^M$ and $0\leq-h<M$.
Every submodule and quotient in \eqref{good:path-tower} is itself
concentrated in $[1-M,0]$. The short exact sequences of the tower
therefore induce conflations in $\H_A^M$.
The direct sums in \eqref{good:path-factors} refine the tower to a
filtration by allowed shifted standards; a zero layer contributes
no factors. Finally,
\[
 0\longrightarrow K_v\longrightarrow P_v
 \longrightarrow\std_v\longrightarrow0
\]
is a short exact sequence of dg modules with all three terms in
$\H_A^M$, so it induces the conflation required in
condition~\ref{ax:dq2}. For $m\geq M$, all the same objects and
factors remain in $\H_A^m$, and every factor shift remains less
than $m$.
\end{proof}

For arrows concentrated in degree zero, $M=1$ and $d_A=0$.
Theorem~\ref{good:every-order} therefore recovers quasi-heredity of
$\Bbbk Q$ for every total order on the vertices.

\begin{example}[Type $A_2$: two orders and two kinds of standards]
\label{good:a2}
Let $A=\Bbbk(1\xrightarrow{a}2)$ with $|a|=-d$, $d\geq1$, and
$d_A=0$. Take $m=d+1$.
Both total orders admit a structure, but their standards differ:
\begin{equation}\label{good:a2-table}
 \begin{array}{c|c|c|c}
 \text{order}&(\std_1,\std_2)&(K_1,K_2)&(\costd_1,\costd_2)\\\hline
 1<_\sigma2&(S_1,S_2)&(S_2[d],0)&(S_1,\dual(Ae_2))\\
 2<_\sigma1&(P_1,S_2)&(0,0)&(S_1,S_2)
 \end{array}
\end{equation}
The costandard entries in \eqref{good:a2-table} are labelled by
vertices; their exceptional sequence is listed in decreasing
order with respect to $\leq_\sigma$.
\begin{enumerate}
 \item For $1<_\sigma2$, the presentation
 \[
  S_2[d]\longrightarrow P_1\longrightarrow S_1
  \longrightarrow S_2[d+1]
 \]
 gives the exceptional sequence $(S_1,S_2)$.
 The shifted factor $S_2[d]$ is necessary when $d\geq1$.
 The costandard $\costd_2=\dual(Ae_2)$ has
 \[
 H^q(\costd_2)\cong
 \begin{cases}S_2,&q=0,\\S_1,&q=d,\\0,&q\notin\{0,d\}.
 \end{cases}
 \]
 Thus the upper endpoint $m-1=d$ of the costandard window is
 attained.
 \item For $2<_\sigma1$, the exceptional sequence is $(S_2,P_1)$.
 The standard $P_1$ has $H^{-d}(P_1)\cong S_2$, so negative
 cohomology in $\std_1$ may contain $S_2$, since $2<_\sigma1$.
 Both kernels vanish.
\end{enumerate}
The costandard calculations follow directly from
$\Hom_{\D_{\fd}(A)}(P_u,S_v[t])\cong H^t(S_v)e_u$ and
$\Hom_{\D_{\fd}(A)}(S_u,\dual(Ae_v)[t])
\cong\dual(H^{-t}(S_u)e_v)$.
\end{example}

\begin{example}[Type $A_4$: standards with negative cohomology]
\label{good:a4}
Let $Q:1\xrightarrow{a}2\xrightarrow{b}3\xrightarrow{c}4$, with
$|a|=|b|=|c|=-1$ and $d_A=0$. Then $M=4$.
Choose $(\sigma(1),\sigma(2),\sigma(3),\sigma(4))=(2,4,1,3)$,
and put $U=P_1/(abA)$. The standards and kernels are
\begin{equation}\label{good:a4-table}
 \begin{array}{c|c|c}
 v&\std_v&K_v\\\hline
 1&U&abA\cong\std_3[2]\\
 2&S_2&bA\cong\std_3[1]\\
 3&P_3&0\\
 4&S_4&0
 \end{array}
 \qquad
 \boldsymbol\std^\sigma=(S_2,S_4,U,P_3).
\end{equation}
Here
\[
 H^{-1}(U)\cong S_2,\qquad H^{-1}(P_3)\cong S_4,
 \qquad 2<_\sigma1,\quad4<_\sigma3.
\]
The kernels $K_1$ and $K_2$ occupy degrees $[-3,-2]$ and
$[-2,-1]$, respectively. Both kernels are single allowed factors in
$\H_A^4$. The cohomology class of $abc$ has degree $-3$, so width
$4$ is already forced by the amplitude of $A$.
\end{example}

\begin{example}[Type $D_4$: a kernel filtered by shifts of two standards]
\label{good:d4}
Let $Q$ have arrows $a:1\to2$, $b:2\to3$, and $c:2\to4$, with
$|a|=|c|=-1$, $|b|=0$, and $d_A=0$. Then $M=3$.
Choose $(\sigma(1),\sigma(2),\sigma(3),\sigma(4))=(1,3,2,4)$,
and put $V=P_2/cA$. We obtain
\begin{equation}\label{good:d4-table}
 \begin{array}{c|c|c}
 v&\std_v&K_v\\\hline
 1&S_1&aA\cong P_2[1]\\
 2&V&cA\cong\std_4[1]\\
 3&S_3&0\\
 4&S_4&0
 \end{array}
 \qquad
 \boldsymbol\std^\sigma=(S_1,S_3,V,S_4).
\end{equation}
The required standard filtration of $K_1$ is
\begin{equation}\label{good:d4-filtration}
 0\subset acA\subset aA,\qquad
 acA\cong\std_4[2],\qquad aA/acA\cong\std_2[1].
\end{equation}
Both factors lie in $\H_A^3$, and their vertex labels satisfy
$1<_\sigma4$ and $1<_\sigma2$. The module $V$ is concentrated in
degree zero, whereas its
shift $V[1]$ occurs as a kernel factor. Thus the projective expression
$K_1\cong P_2[1]$ must be refined as in
\eqref{good:d4-filtration} to verify condition~\ref{ax:dq2}.
\end{example}

\begin{example}[A non-trivial differential and an excluded width]
\label{good:differential}
Fix $d\geq0$. Let $Q$ have arrows $a:1\to2$, $b:2\to3$, and
$c:1\to3$, and define
\begin{equation}\label{good:differential-algebra}
 |a|=0,\quad |b|=-d,\quad |c|=-d-1,\qquad
 d_A(c)=ab,\quad d_A(a)=d_A(b)=0.
\end{equation}
The path-degree bound is $M=d+2$. The pair
$c\xrightarrow{d_A}ab$ is acyclic, so
\begin{equation}\label{good:differential-cohomology}
 H^*(A)\cong\Bbbk(1\xrightarrow{a}2\xrightarrow{b}3)/(ab)
 \quad\text{as graded algebras},\qquad |b|=-d.
\end{equation}
The least integer $m$ with $A\in\H_A^m$ is therefore $d+1$.
For the order $1<_\sigma3<_\sigma2$, the objects in
\eqref{good:path-standards} are
\begin{equation}\label{good:differential-standards}
 (\std_1,\std_2,\std_3)=(S_1,P_2,S_3),\qquad
 K_1=\Bbbk\{a,ab,c\},\qquad K_2=K_3=0.
\end{equation}
The only required kernel filtration is
\begin{equation}\label{good:differential-filtration}
 0\subset aA\subset K_1,\qquad
 aA\cong\std_2,\qquad K_1/aA\cong\std_3[d+1].
\end{equation}
Consequently the order $1<_\sigma3<_\sigma2$ admits a structure at
width $d+2$. Nevertheless, $K_1\simeq S_2$, and all the projectives,
standards, and kernels in \eqref{good:differential-standards} already
belong to $\H_A^{d+1}$.

The failure at width $d+1$ is detected by the costandard
$\costd_3$. The triangles
\begin{equation}\label{good:differential-triangles}
 \begin{gathered}
 S_3[d]\longrightarrow P_2\longrightarrow S_2
 \longrightarrow S_3[d+1],\\
 S_2\longrightarrow P_1\longrightarrow S_1\longrightarrow S_2[1]
 \end{gathered}
\end{equation}
and the standard--costandard pairing give, for every $t\in\ZZ$,
\begin{equation}\label{good:differential-costandard}
 \begin{aligned}
 H^t(\costd_3)e_1
 &=\Hom_{\D_{\fd}(A)}(P_1,\costd_3[t])\\
 &\cong\Hom_{\D_{\fd}(A)}(S_2,\costd_3[t])\\
 &\cong\Hom_{\D_{\fd}(A)}(S_3[d],\costd_3[t-1])\\
 &\cong
 \begin{cases}\Bbbk,&t=d+1,\\0,&t\neq d+1.
 \end{cases}
 \end{aligned}
\end{equation}
In addition, $H^t(\costd_3)e_2=0$ for every $t$, and
$H^t(\costd_3)e_3\cong\Bbbk$ precisely when $t=0$.
Proposition~\ref{adv:canonical} determines the standard objects
from $(A,\sigma)$, independently of the width.
Hence the costandard window in
Proposition~\ref{adv:costandards} and
\eqref{good:differential-costandard} force $m\geq d+2$ for every
structure with the order $1<_\sigma3<_\sigma2$.
The least width is exactly $d+2$.

For $d=0$, the algebra in \eqref{good:differential-cohomology} is
ordinary, and the standard objects form a full exceptional sequence
in $\D^b(\mod H^0(A))$. Example~\ref{bad:exceptional-filtration}
gives an elementary proof that $(S_1,P_2,S_3)$ has no kernel filtration
as required by condition~\ref{ax:dq2} at width $1$.
For $d\geq1$, the same obstruction occurs with genuine negative
cohomology in $A$. Thus cohomological amplitude alone does not
determine the width of an $m$-qh structure, even for a smooth dg
path algebra.
\end{example}

\section{Excluded non-smooth algebras}\label{bad:nonsmooth}

The algebras in Examples~\ref{bad:negative-diagonal},
\ref{bad:without-filtration}, \ref{bad:ordinary-diagonal}, and
\ref{bad:cycle} are not homologically smooth. Excluding those algebras is a
consequence of conditions~\ref{ax:dq1} and~\ref{ax:dq2}, and is a
required feature of the definition.

\begin{lemma}[Resolutions over graded dual numbers]\label{bad:dualnumbers-resolution}
For an integer $\ell\geq0$, set
\begin{equation}\label{bad:dualnumbers-algebra}
 R_\ell=\Bbbk\oplus\Bbbk\varepsilon,
 \qquad\varepsilon^2=0,\qquad |\varepsilon|=-\ell,
 \qquad d_{R_\ell}=0,\qquad S=R_\ell/(\varepsilon).
\end{equation}
Then
\begin{equation}\label{bad:dualnumbers-ext}
 \Hom_{\D_{\fd}(R_\ell)}(S,S[t])\cong
 \begin{cases}
 \Bbbk,&t\in(\ell+1)\ZZ_{\geq0},\\
 0,&\text{otherwise}.
 \end{cases}
\end{equation}
In particular, $S\notin\per R_\ell$ and $R_\ell$ is not
homologically smooth.
\end{lemma}
\begin{proof}
A semifree resolution $F\to S$ has
\begin{equation}\label{bad:dualnumbers-semifree}
 F=\bigoplus_{a\geq0}x_aR_\ell,\qquad
 |x_a|=-a(\ell+1),\qquad
 d_F(x_0)=0,\qquad d_F(x_a)=x_{a-1}\varepsilon\quad(a\geq1).
\end{equation}
The augmentation $F\to S$ sends $x_0$ to $1$ and every $x_a$ with
$a>0$ to zero. The relation $\varepsilon^2=0$ gives $d_F^2=0$.
Each $x_a\varepsilon$ is the boundary of $x_{a+1}$, so the
augmentation is a quasi-isomorphism. Applying the dg Hom complex
from $F$ to $S$ gives zero differential and one basis vector in each
degree $a(\ell+1)$, proving \eqref{bad:dualnumbers-ext}.
For a perfect object $X$, the total graded vector space
$\bigoplus_t\Hom_{\D_{\fd}(R_\ell)}(X,S[t])$ is finite-dimensional.
Thus \eqref{bad:dualnumbers-ext} proves $S\notin\per R_\ell$.
Smoothness would imply that every object of $\D_{\fd}(R_\ell)$ is
perfect: tensor a finite construction of the diagonal bimodule
with the object. Hence $R_\ell$ is not smooth.
\end{proof}

\begin{example}[Negative cohomology at the standard's own vertex]
\label{bad:negative-diagonal}
Fix $\ell\geq1$, let $A=R_\ell$ be the algebra
\eqref{bad:dualnumbers-algebra}, and take $m=\ell+1$. With one
vertex, choose
\[
 \std_1=P_1=R_\ell,\qquad K_1=0.
\]
Condition~\ref{ax:dq2} holds. Every support restriction at a strictly
larger vertex is vacuous, and
\[
 \End_{\H_{R_\ell}^{\ell+1}}(\std_1)=\Bbbk,
 \qquad [H^0(\std_1):S]=1.
\]
Nevertheless,
\begin{equation}\label{bad:negative-diagonal-self}
 H^{-\ell}(\std_1)e_1=\Bbbk\varepsilon,
 \qquad
 \Hom_{\D_{\fd}(R_\ell)}(\std_1,\std_1[-\ell])=\Bbbk.
\end{equation}
Thus allowing the standard's own simple in negative cohomology
admits a non-smooth algebra, even when the standard is a brick and
satisfies condition~\ref{ax:dq2}. The vanishing of the first vector
space in \eqref{bad:negative-diagonal-self} is essential.
More generally, for one vertex,
condition~\ref{ax:dq2} forces $\std_1\cong A$, and
condition~\ref{ax:dq1} then forces $H^*(A)=\Bbbk$ in degree zero.
Increasing $m$ never makes $R_\ell$ satisfy both axioms.
\end{example}

\begin{example}[Cohomological support does not replace the filtration]
\label{bad:without-filtration}
Take $A=R_1$, $m=2$, and $\std_1=S=A/(\varepsilon)$.
The object $\std_1$ satisfies condition~\ref{ax:dq1}. In particular,
$\std_1$ is a brick. The short exact sequence
\begin{equation}\label{bad:simple-presentation}
 0\longrightarrow S[1]\xrightarrow{\iota}A
 \xrightarrow{p}S\longrightarrow0,
 \qquad\iota(1)=\varepsilon,
\end{equation}
induces a conflation in $\H_A^2$. There is no vertex $w>_\sigma1$,
so condition~\ref{ax:dq2} requires $K_1=0$, whereas the
kernel in \eqref{bad:simple-presentation} is $S[1]\neq0$.
If shifts of $\std_1$ were admitted as factors of $K_1$, then
$K_1=\std_1[1]$ would
satisfy the weakened filtration requirement.
Lemma~\ref{bad:dualnumbers-resolution} gives
\[
 \std_1\notin\per A,
 \qquad
 \Hom_{\D_{\fd}(A)}(\std_1,\std_1[2])=\Bbbk.
\]
Hence even condition~\ref{ax:dq1} and an ordinary dg quotient
$P_1\twoheadrightarrow\std_1$ do not imply perfection,
exceptionality, or smoothness. Requiring factors $\std_w[s]$ of
$K_v$ to satisfy $w>_\sigma v$ in condition~\ref{ax:dq2} is indispensable.
\end{example}

\begin{example}[The multiplicity-one condition in degree zero]
\label{bad:ordinary-diagonal}
Let $A=R_0=\Bbbk[x]/(x^2)$ be concentrated in degree zero, take
$m=1$, and set $\std_1=P_1=A$, $K_1=0$. Then
condition~\ref{ax:dq2} holds and every off-diagonal support or
orthogonality condition is vacuous. However,
\[
 [\std_1:S_1]=2,\qquad \End_A(\std_1)=A\not\cong\Bbbk.
\]
Lemma~\ref{bad:dualnumbers-resolution} proves that $A$ is not smooth.
Thus the requirement $[H^0(\std_v):S_v]=1$ in
condition~\ref{ax:dq1} is necessary
already for $m=1$; off-diagonal projective tests and the kernel
filtration alone do not recover classical quasi-heredity.
\end{example}

\begin{example}[Exceptional terms and backward degree-zero Hom are insufficient]
\label{bad:cycle}
Let $A$ be the dg algebra with arrows
$a:1\to2$ and $b:2\to1$, relations $ab=ba=0$, degrees
$|a|=|b|=-1$, and zero differential. Paths are concatenated from left
to right. Set $m=2$, order $1<_\sigma2$, and
\[
 \std_1=P_1,\qquad\std_2=P_2,\qquad K_1=K_2=0.
\]
Condition~\ref{ax:dq2} holds. Since $e_iAe_i=\Bbbk$, each
$\std_i$ is exceptional. Moreover,
\begin{equation}\label{bad:cycle-backward}
 \Hom_{\H_A^2}(\std_2,\std_1)=0,
 \qquad
 \Hom_{\D_{\fd}(A)}(\std_2,\std_1[-1])=\Bbbk a.
\end{equation}
The cohomological support restriction in condition~\ref{ax:dq1}
excludes $H^{-1}(\std_1)e_2=\Bbbk a$.

To verify non-smoothness directly, let $F\to S_1$ be the semifree
resolution with a generator $x_j$ in degree $-2j$, at vertex $1$
for even $j$ and vertex $2$ for odd $j$. Define
\begin{equation}\label{bad:cycle-resolution}
 d_F(x_0)=0,\qquad
 d_F(x_j)=
 \begin{cases}
 x_{j-1}a,&j\text{ odd},\\
 x_{j-1}b,&j>0\text{ even}.
 \end{cases}
\end{equation}
The augmentation sends $x_0$ to a generator of $S_1$ and every
$x_j$ with $j>0$ to zero. The relations give $d_F^2=0$; every radical
basis vector is the boundary of the next generator. Thus $F\to S_1$
is a quasi-isomorphism, and
\begin{equation}\label{bad:cycle-ext}
 \Hom_{\D_{\fd}(A)}(S_1,S_1[t])\cong
 \begin{cases}\Bbbk,&t\in4\ZZ_{\geq0},\\0,&\text{otherwise}.
 \end{cases}
\end{equation}
Therefore $S_1\notin\per A$ and $A$ is not smooth. In particular,
adding individual exceptionality to degree-zero backward
orthogonality still does not protect smoothness.
\end{example}

\section{Rejected standards and orders on smooth algebras}
\label{bad:smooth}

Failure of chosen standards to satisfy the axioms need not indicate
a singularity of the algebra. Examples~\ref{bad:fork} and
\ref{bad:forbidden-order} distinguish a rejected choice of standards
from an order for which no choice can satisfy the axioms.

\begin{example}[Positive shifts of the target detect only $H^0$]
\label{bad:fork}
For $d\in\{1,2\}$, let
\begin{equation}\label{bad:fork-algebra}
 A_d=\Bbbk(1\xrightarrow{a}3\xleftarrow{b}2),\qquad
 |a|=-d,\quad |b|=-1,\quad d_{A_d}=0,
 \qquad m=d+1,
\end{equation}
and use the order $1<_\sigma2<_\sigma3$. Consider
\begin{equation}\label{bad:fork-standards}
 (\std_1,\std_2,\std_3)=(P_1,S_2,P_3),
 \qquad (K_1,K_2,K_3)=(0,P_3[1],0).
\end{equation}
The only non-identity presentation is induced by
\begin{equation}\label{bad:fork-presentation}
 0\longrightarrow P_3[1]\xrightarrow{\iota_b}P_2
 \longrightarrow S_2\longrightarrow0,
 \qquad\iota_b(e_3)=b.
\end{equation}
All three objects in \eqref{bad:fork-presentation} belong to
$\H_{A_d}^{d+1}$, and $K_2=\std_3[1]$ is an allowed factor.
Thus condition~\ref{ax:dq2} holds. Each $\std_i$ is exceptional,
and $H^0(\std_i)=S_i$. Consequently,
\begin{equation}\label{bad:fork-positive-tests}
 \Hom_{\D_{\fd}(A_d)}(P_w,\std_v[k])
 =H^k(\std_v)e_w=0
 \qquad(v<_\sigma w,\ 0\leq k<m).
\end{equation}
For any standard in $\H_A^m$, every test with $k>0$ in
\eqref{bad:fork-positive-tests} vanishes automatically.
Writing $\Hom_{\H_A^m}(P_w,\std_v[k])$ instead does not change the
calculation when $\std_v[k]\in\H_A^m$; the latter Hom space is
undefined when $\std_v[k]\notin\H_A^m$.

The actual backward morphism spaces for
\eqref{bad:fork-standards} are
\begin{equation}\label{bad:fork-homs}
 \begin{aligned}
 \Hom_{\D_{\fd}(A_d)}(\std_2,\std_1[t])
 &=\begin{cases}\Bbbk,&t=2-d,\\0,&t\neq2-d,\end{cases}\\
 \Hom_{\D_{\fd}(A_d)}(\std_3,\std_1[t])
 &=\begin{cases}\Bbbk,&t=-d,\\0,&t\neq-d,\end{cases}\\
 \Hom_{\D_{\fd}(A_d)}(\std_3,\std_2[t])&=0.
 \end{aligned}
\end{equation}
Indeed, $\Hom_{\D_{\fd}(A_d)}(P_2,P_1[t])=0$ for all $t$, so
\eqref{bad:fork-presentation} gives
\[
 \Hom_{\D_{\fd}(A_d)}(S_2,P_1[t])
 \cong\Hom_{\D_{\fd}(A_d)}(P_3[1],P_1[t-1])
 \cong H^{t-2}(P_1)e_3.
\]
The other two equalities in \eqref{bad:fork-homs} are projective
calculations. Thus:
\begin{enumerate}[label=\textup{(\roman*)}]
 \item for $d=1$, all backward degree-zero Hom spaces vanish,
 but $\Hom_{\D_{\fd}(A_1)}(\std_2,\std_1[1])=\Bbbk$;
 \item for $d=2$, even backward degree-zero orthogonality fails:
 $\Hom_{\H_{A_2}^3}(\std_2,\std_1)=\Bbbk$.
\end{enumerate}
For both values of $d$, condition~\ref{ax:dq1} rejects the chosen
$\std_1=P_1$, since $H^{-d}(P_1)e_3=\Bbbk a$.
The algebra $A_d$ does admit an $m$-qh structure for the same order:
choose $\std'_i=S_i$ for every $i$, with
\[
 K'_1=P_3[d],\qquad K'_2=P_3[1],\qquad K'_3=0.
\]
Conditions~\ref{ax:dq1} and~\ref{ax:dq2} hold for
$(\std'_1,\std'_2,\std'_3)$, so $A_d$ is smooth. The failure in
\eqref{bad:fork-standards} concerns the chosen standards, rather
than the existence of an $m$-qh structure.
\end{example}

\begin{example}[A smooth algebra with an order excluded at every width]
\label{bad:forbidden-order}
Let $A$ be the ordinary algebra with arrows $a:1\to2$ and
$b:2\to1$, and relation $ab=0$. A basis of $A$ is
$e_1,e_2,a,b,ba$, with paths concatenated from left to right.
The right projectives satisfy
\begin{equation}\label{bad:auslander-radicals}
 \rad P_1=aA\cong S_2,\qquad
 \rad P_2=bA\cong P_1.
\end{equation}
Thus $\operatorname{pdim}_A S_2=1$ and
$\operatorname{pdim}_A S_1=2$, whence
$\operatorname{gldim}A=2$. In particular, $A$ is homologically smooth.

Take the order $1<_\sigma2$. For every $m\geq1$,
the greatest vertex is $2$, so condition~\ref{ax:dq2} forces
$\std_2\cong P_2$. But
\begin{equation}\label{bad:auslander-diagonal}
 H^0(P_2)e_2=e_2Ae_2=\Bbbk e_2\oplus\Bbbk ba,
\end{equation}
contradicting $\dim_{\Bbbk}H^0(\std_2)e_2=1$ in
condition~\ref{ax:dq1}. Therefore no $m$-qh structure exists for
the order $1<_\sigma2$, at any width.

For the opposite order $2<_\sigma1$, choose
\[
 \std_1=P_1,\qquad\std_2=S_2,\qquad
 K_1=0,\qquad K_2=bA\cong\std_1.
\]
Equations~\eqref{bad:auslander-radicals} show that both axioms hold
already for $m=1$. Hence increasing the width does not remove every
obstruction arising from the choice of order.
\end{example}

\begin{example}[Full exceptionality does not imply the kernel filtration]
\label{bad:exceptional-filtration}
Consider the ordinary algebra
\begin{equation}\label{bad:relation-algebra}
 B=\Bbbk(1\xrightarrow{a}2\xrightarrow{b}3)/(ab),
 \qquad 1<_\sigma3<_\sigma2,
\end{equation}
and the quotient modules
\begin{equation}\label{bad:relation-standards}
 (\std_1,\std_2,\std_3)=(S_1,P_2,S_3),
 \qquad(K_1,K_2,K_3)=(S_2,0,0).
\end{equation}
The standards satisfy condition~\ref{ax:dq1}. The sequence
$(\std_1,\std_3,\std_2)=(S_1,S_3,P_2)$ is a full exceptional
sequence in $\D_{\fd}(B)$:
\begin{enumerate}[label=\textup{(\roman*)}]
 \item the projectives $S_3=P_3$ and $P_2$ are exceptional;
 \item the resolution
 $0\to P_3\to P_2\to P_1\to S_1\to0$ makes $S_1$ exceptional;
 \item all backward derived Hom spaces vanish, since their
 sources are $P_2$ or $P_3$, and the corresponding vertex components
 of their targets vanish;
 \item the triangle $S_3\to P_2\to S_2\to S_3[1]$ puts $S_2$ in
 the thick subcategory generated by $S_3,P_2$, so the standards
 generate every simple and hence $\D_{\fd}(B)$.
\end{enumerate}
Nevertheless, condition~\ref{ax:dq2} fails for $m=1$. A filtration
of $K_1=S_2$ in $\mod B$ by $\std_2=P_2$ and $\std_3=S_3$,
whose labels satisfy $2>_\sigma1$ and $3>_\sigma1$, would give
\begin{equation}\label{bad:relation-multiplicities}
 [S_2]=x[P_2]+y[S_3]=x[S_2]+(x+y)[S_3],
 \qquad x,y\in\ZZ_{\geq0},
\end{equation}
in $K_0(\mod B)$, forcing $x=1$, $y=-1$. Thus even a full
exceptional sequence satisfying condition~\ref{ax:dq1}, together
with quotient presentations $P_v\twoheadrightarrow\std_v$, does
not imply condition~\ref{ax:dq2}. In fact, no width-$1$ structure
exists for the order in \eqref{bad:relation-algebra}:
condition~\ref{ax:dq2} forces $\std_2=P_2$, while
condition~\ref{ax:dq1} and the quotient presentations force
$\std_3=S_3$ and $\std_1=S_1$.

For $m=2$, the same ordinary algebra and the same standards do
satisfy condition~\ref{ax:dq2}. Rotate the short exact sequence
$0\to S_3\to P_2\to S_2\to0$ to obtain the conflation
\begin{equation}\label{bad:relation-shifted-filtration}
 P_2\longrightarrow S_2\longrightarrow S_3[1]\dashrightarrow
 \qquad\text{in }\H_B^2.
\end{equation}
The two-factor filtration
$0\to P_2\to S_2$ has factors $\std_2$ and $\std_3[1]$.
The vertex labels of both factors satisfy $2>_\sigma1$ and
$3>_\sigma1$, so
$K_1\in\Filt_{\H_B^2}\{\std_2,\std_3[1]\}$.
The map $P_2\to S_2$ in \eqref{bad:relation-shifted-filtration}
is an inflation in $\H_B^2$, although it is a surjection of
ordinary modules. Thus an extriangulated filtration must not be
replaced by a chain of dg submodules.
Example~\ref{good:differential} gives a path dga with non-trivial
differential realising the same distinction between widths $1$ and $2$.
\end{example}

\section{Summary of consequences and counterexamples}\label{sec:summary}

Conditions~\ref{ax:dq1} and \ref{ax:dq2} have the consequences in
Table~\ref{tab:consequences}. The negative examples in
Table~\ref{tab:counterexamples} identify failed weakenings of the
axioms and genuine restrictions on orders and widths.

\begin{table}[htbp]
\centering\small
\begin{tabularx}{\linewidth}{@{}>{\raggedright\arraybackslash}p{.42\linewidth}>{\raggedright\arraybackslash}X@{}}
\toprule
Consequence of (dq1) and (dq2) & Precise result\\\midrule
Every standard is a brick and exceptional & Theorem~\ref{thm:equivalence}.\\
Backward orthogonality in every degree & Equation~\eqref{eq:backward}.\\
Standards are full in $\D_{\fd}(A)=\per A$ & Theorems~\ref{thm:smoothness} and \ref{thm:equivalence}.\\
Homological smoothness & A theorem, without an additional axiom;
Theorem~\ref{thm:smoothness}.\\
Canonical standard objects & Determined by $(A,\sigma)$;
Proposition~\ref{adv:canonical}.\\
Costandards $\costd_v\in\H_A^m[1-m]$ & Proposition~\ref{adv:costandards}.\\
Recognition of allowed filtrations & Proposition~\ref{adv:criterion}.\\
Least width and all larger widths & Corollary~\ref{adv:width}.\\
Classical quasi-heredity at $m=1$ & Corollary~\ref{cor:classical}.\\
Injective presentations and passage to $A^{\op}$ & Proposition~\ref{adv:opposite}.\\
Euler Cartan determinant $1$ & Proposition~\ref{adv:kzero}.\\
All vertex orders for graded acyclic path dgas & At the path-degree
bound; Theorem~\ref{good:every-order}.\\
\bottomrule
\end{tabularx}
\caption{Properties proved from the two axioms.}
\label{tab:consequences}
\end{table}

\begin{table}[htbp]
\centering\small
\begin{tabularx}{\linewidth}{@{}>{\raggedright\arraybackslash}p{.42\linewidth}>{\raggedright\arraybackslash}X@{}}
\toprule
Claim disproved & Counterexample and reason\\\midrule
The restrictions $H^q(\std_v)e_v=0$ for $q<0$ can be omitted &
Example~\ref{bad:negative-diagonal}: graded dual numbers are
non-smooth despite a brick projective standard.\\
Composition factors and a quotient presentation suffice &
Example~\ref{bad:without-filtration}: the simple over dg dual
numbers satisfies (dq1), but its kernel is a shift of the same
standard.\\
The requirement $[H^0(\std_v):S_v]=1$ is dispensable &
Example~\ref{bad:ordinary-diagonal}: ordinary dual numbers.\\
Exceptional terms and backward degree-zero Hom protect smoothness &
Example~\ref{bad:cycle}: a graded radical-square-zero cycle.\\
Positive target shifts test the required support &
Example~\ref{bad:fork}: the tests miss backward morphisms in
degree $1$ or degree $0$.\\
Smoothness permits every vertex order at some width &
Example~\ref{bad:forbidden-order}: the greatest vertex is $2$, but
$\dim_{\Bbbk}H^0(P_2)e_2=2$.\\
A full exceptional sequence implies (dq2) &
Example~\ref{bad:exceptional-filtration}: a required classical
filtration would have a negative multiplicity.\\
The least $m$ with $A\in\H_A^m$ suffices &
Example~\ref{good:differential}: a required costandard degree
forces a larger width.\\
All filtrations have invariant unsigned multiplicities &
Example~\ref{adv:zero-filtration}: zero has a two-factor filtration,
even over a semisimple algebra.\\
Extriangulated filtrations are chains of submodules &
Equation~\eqref{bad:relation-shifted-filtration}: a surjection of
modules is an inflation in $\H_B^2$.\\
\bottomrule
\end{tabularx}
\caption{Failures of weakened axioms or stronger classical expectations.
No non-smooth algebra satisfies both (dq1) and (dq2).}
\label{tab:counterexamples}
\end{table}

\clearpage

\begin{thebibliography}{RS22}
\bibitem[BB26]{BB26} A. Bodzenta and A. Bondal,
\emph{Highest weight categories via pairs of dual exceptional sequences},
2026. \href{https://arxiv.org/abs/2601.22004}{arXiv:2601.22004}.

\bibitem[MP25]{MP25} N. Mochizuki and M. Plogmann,
\emph{On the Auslander--Reiten theory for extended hearts of proper
connective dg algebras}, 2025.
\href{https://arxiv.org/abs/2505.16560}{arXiv:2505.16560}.

\bibitem[RS22]{RS22} T. Raedschelders and G. Stevenson,
\emph{Proper connective differential graded algebras and their geometric
realizations}, European J. Math. \textbf{8} (2022), S574--S598.
\href{https://arxiv.org/abs/1903.02849}{arXiv:1903.02849}.

\end{thebibliography}
\end{document}
